\begin{onehalfspace}

Il Capitolo~\ref{chap:obiettivi} ha fissato l'ipotesi da verificare: ridurre il numero
medio di esecuzioni necessarie per ottenere fattori validi di Shor mediante un filtro
classico. Il Capitolo~\ref{chap:quadro} ha poi definito rumore, metriche e metodologia.
Qui si presenta soltanto l'architettura che collega tali elementi alle due campagne
sperimentali: la prima confronta TOP-1, TOP-4 e il filtro \ac{ML}; la seconda studia
codici correttori e decodifica della sindrome. Le motivazioni estese degli strumenti e
gli estratti del codice sono raccolti, rispettivamente, nelle
Sezioni~\ref{app:strumenti} e~\ref{app:codice}; i risultati sono nel
Capitolo~\ref{chap:risultati}.

L'architettura ha quattro livelli. Il primo costruisce e compila il circuito, una volta per configurazione. Il secondo applica il rumore e genera i campioni. Il terzo esegue il post-processing, il classificatore o il decoder. Il quarto salva metriche e metadati e produce tabelle e figure. Così, cambiare il classificatore non modifica il circuito e una compilazione diversa non viene scambiata per un beneficio del metodo. Nella prima fase le strategie condividono l'istogramma delle misure; nella seconda, la sequenza dei detection event. In entrambi i casi il confronto usa gli stessi campioni.

\section{Strumenti e configurazione dell'ambiente}
\label{chap:strumenti}

\subsection{Strumenti e ambiente di simulazione}
\label{sec:strumenti_ambiente}

Gli strumenti sono scelti in base ai circuiti. \textbf{Qiskit}~\cite{qiskit2023} e \textbf{Qiskit Aer} costruiscono, compilano e simulano Shor sulle piccole istanze. Il classificatore usa \textbf{scikit-learn}~\cite{scikit_learn}. Per i milioni di campioni del surface code si usano \textbf{Stim}~\cite{stim2021} e \textbf{PyMatching}~\cite{pymatching2021}. Stim è efficiente sui circuiti stabilizer, mentre una simulazione statevector richiede uno spazio che cresce come $2^n$. Questa efficienza non si estende a Shor, che contiene rotazioni non-Clifford. Il confronto fra framework e metodi di Aer è nella Sezione~\ref{app:strumenti}.

La campagna rumorosa principale usa
\texttt{AerSimulator(method='statevector')} per il circuito didattico a 12 qubit di
$N=15$; $N=21$ e $N=35$ sono riservati ai controlli ideali di correttezza e costo.
Il noise model uniforme e illustrativo combina depolarizzazione, rilassamento termico e
readout e viene applicato, dopo compilazione, alla base \texttt{rz/sx/x/cx}. Non è una
calibrazione di una QPU nominata.

\subsection{Ambiente e contratto operativo}
\label{sec:setup}

Le esecuzioni canoniche avvengono su \ac{CPU} in \ac{WSL}2 con dipendenze fissate;
interprete, librerie e backend sono registrati nel manifest di ciascun artefatto.
L'accelerazione \ac{GPU}, prevista inizialmente, non è entrata nel contratto
riproducibile perché il backend Aer non era utilizzabile nell'ambiente di sviluppo.
\label{sec:gpu_limit} Versioni, configurazione WSL, limite di memoria e diagnosi del
fallimento sono documentati nelle Sezioni~\ref{sec:wsl} e~\ref{sec:gpu}; il comando di
installazione e l'errore osservato compaiono nei
Listati~\ref{lst:installazione_dellambiente_python_e}
e~\ref{lst:errore_gpu_su_wsl}.

\end{onehalfspace}
